\section*{Chapitre \ref{ch:b1:systems-of-points} --- Systèmes de points~; introduction aux solides}

\begin{solution}{exo:b1:systems-of-points:1}
$x_G = (2.0 \times 0 + 3.0 \times 1.0)/5.0 = \qty{0.60}{m}$. Terre--Lune~:
$384000/82 = \qty{4700}{km}$ du centre de la Terre, donc à l'intérieur de
la Terre (\qty{6370}{km})~: la Terre oscille autour d'un point situé
\qty{1700}{km} sous sa surface.
\end{solution}

\begin{solution}{exo:b1:systems-of-points:2}
$4.0v = 0.010 \times 800$~: $v = \qty{2.0}{m/s}$. Balle $\tfrac12 \times 0.010
\times 800^2 = \qty{3.2}{kJ}$, fusil \qty{8}{J}~: cette énergie vient de
la poudre~; le corps léger en prend presque tout ($E_k = p^2/2m$ à $p$ égal).
\end{solution}

\begin{solution}{exo:b1:systems-of-points:3}
$v = 20000/2500 = \qty{8.0}{m/s}$~; $E_k$~: \qty{200}{kJ} avant, $\tfrac12 \times
2500 \times 64 = \qty{80}{kJ}$ après~: $60\%$ perdus.
\end{solution}

\begin{solution}{exo:b1:systems-of-points:4}
Extrémité~: $\int_0^L x^2(M/L)\dd x = ML^2/3$~; centre~: $\int_{-L/2}^{L/2}x^2(M/L)
\dd x = ML^2/12$~; Huygens avec $d = L/2$~: $ML^2/12 + ML^2/4 = ML^2/3$.
\end{solution}

\begin{solution}{exo:b1:systems-of-points:5}
$E_k = \tfrac12 Mv^2(1 + k)$~: cylindre $1.5$, cerceau $2$ fois $\tfrac12 Mv^2$.
À vitesse égale, le cerceau a plus d'énergie~: il roule plus loin sur un
sol plat (à frottement égal) et monte plus haut ($h = v^2(1 + k)/2g$).
\end{solution}

\begin{solution}{exo:b1:systems-of-points:6}
HCl~: $\mu = 35m_H/36 = 0.972m_H$~; atome~: $\mu = m_e/(1 + 1/1836) =
0.99946m_e$. Les fréquences varient en $\propto 1/\sqrt\mu$~: $+1.4\%$ pour HCl, $+0.027\%$
pour l'hydrogène (le décalage isotopique entre les spectres de H et de D en dépend).
\end{solution}

\begin{solution}{exo:b1:systems-of-points:7}
$J = ML^2/3$, $d = L/2$~: $T = 2\pi\sqrt{(ML^2/3)/(MgL/2)} = 2\pi\sqrt{2L/3g}
= \qty{1.64}{s}$~; longueur équivalente $2L/3 = \qty{0.67}{m}$. Un coup au
centre de percussion met la tige en rotation autour du pivot sans
réaction impulsive en ce point~: frapper la balle au «~point idéal~» de
la batte, et les mains ne ressentent rien.
\end{solution}

\begin{solution}{exo:b1:systems-of-points:8}
$J = \tfrac12 MR^2 = \qty{1.0e-3}{kg.m^2}$~; $C = 4\pi^2J/T^2 = \qty{9.9e-3}{N.m/rad}$.
\end{solution}

\begin{solution}{exo:b1:systems-of-points:9}
$\omega = \qty{314}{rad/s}$, $E = \tfrac12 J\omega^2 = \qty{4.9e5}{J}$~;
$\Gamma = J\omega/t = \qty{314}{N.m}$~; tours~: $\tfrac12\omega t/2\pi = 250$.
\end{solution}

\begin{solution}{exo:b1:systems-of-points:10}
$a = g\sin\alpha/(1 + k)$~: sphère ($k = 2/5$) $0.714$, cylindre ($1/2$)
$0.667$, cerceau ($1$) $0.500$ fois $g\sin\alpha$~: la sphère d'abord, le
cerceau en dernier~; ni masse ni rayon nulle part.
\end{solution}

\begin{solution}{exo:b1:systems-of-points:11}
Quantité de mouvement~: $0.010 \times 400 = 2.01\,v$~: $v = \qty{2.0}{m/s}$~;
puis l'énergie~: $h = v^2/2g = \qty{0.20}{m}$. $E_k$~: \qty{800}{J} avant, $\tfrac12 \times
2.01 \times 4.0 = \qty{4.0}{J}$ après~: $99.5\%$ perdus en déformation et
en chaleur. Le choc est inélastique --- l'énergie ne s'y conserve pas,
la quantité de mouvement si (la tension du fil est verticale et le choc bref).
\end{solution}

\begin{solution}{exo:b1:systems-of-points:12}
$m_1v_1 = m_1v_1' + m_2v_2'$ et $m_1v_1^2 = m_1v_1'^2 + m_2v_2'^2$ donnent
$v_1 + v_1' = v_2'$ (en divisant l'équation d'énergie par celle de la
quantité de mouvement), d'où les formules. Fraction d'énergie transmise à la cible $4m_1m_2/(m_1 + m_2)^2$~:
proton $100\%$~; carbone $4 \times 12/169 = 28\%$~; plomb $4 \times 207/208^2 =
1.9\%$. Les noyaux légers prennent l'énergie du neutron en quelques chocs
(eau, graphite)~; il en faudrait des centaines avec le plomb.
\end{solution}

\begin{solution}{pb:b1:systems-of-points:1}
\textbf{1.} Sur l'axe, au centre du disque, par symétrie.

\textbf{2.} König~: $E_k = \tfrac12 Mv_G^2 + E_k^*$, et dans le \omterm{def:b1:systems-of-points:barycentric}{référentiel
barycentrique} le yo-yo tourne autour de son axe~: $E_k^* = \tfrac12 J\omega^2$.

\textbf{3.} Anneau de rayon $\rho$ et de masse $2M\rho\,\dd\rho/R^2$~: $J = \int_0^R
\rho^2\cdot 2M\rho\,\dd\rho/R^2 = MR^2/2 = 0.5 \times 0.050 \times 9\times10^{-4}
= \qty{2.25e-5}{kg.m^2}$.

\textbf{4.} La ficelle est fixe et se déroule de l'axe~: le point de
l'axe en contact avec elle est momentanément immobile, donc $v_G = r\omega$.

\textbf{5.} $E_k = \tfrac12 Mv_G^2 + \tfrac12(\tfrac12 MR^2)(v_G/r)^2 = \tfrac12
Mv_G^2(1 + R^2/2r^2)$~; $R^2/2r^2 = 9/(2 \times 0.16) = 28$~: facteur $29$ ---
presque toute l'énergie est dans la rotation.

\textbf{6.} Poids $Mg$ vers le bas, tension $T$ vers le haut~: $Ma_G = Mg - T$.

\textbf{7.} Seule la tension a un moment par rapport à l'axe (\omterm{def:b1:angular-momentum:moment}{bras de levier} $r$)~:
$J\dot\omega = Tr$.

\textbf{8.} $\dot\omega = a_G/r$~: $T = Ja_G/r^2 = (R^2/2r^2)Ma_G$~; puis
$Ma_G = Mg - (R^2/2r^2)Ma_G$, d'où $a_G = g/29 = \qty{0.34}{m/s^2}$.

\textbf{9.} $T = M(g - a_G) = 0.050 \times 9.47 = \qty{0.47}{N}$, soit $97\%$
du poids~: la ficelle le retient presque entièrement.

\textbf{10.} $t = \sqrt{2h/a_G} = \qty{2.4}{s}$~; $v_G = a_Gt = \qty{0.82}{m/s}$.
Une pierre~: \qty{0.45}{s} et \qty{4.4}{m/s}.

\textbf{11.} $\omega = v_G/r = \qty{205}{rad/s} = \qty{1960}{rpm}$.

\textbf{12.} $\tfrac12 Mv_G^2 = \qty{0.017}{J}$ contre $\tfrac12 J\omega^2 =
\qty{0.47}{J}$~: $28$ fois plus d'énergie dans la rotation que dans la chute.

\textbf{13.} $Mgh = \qty{0.49}{J}$~; $\tfrac12 Mv_G^2 \times 29 = 0.5 \times 0.050
\times 0.67 \times 29 = \qty{0.49}{J}$. La tension s'applique au point de
la ficelle qui est immobile (la ficelle ne bouge pas)~: puissance nulle.

\textbf{14.} La force d'arrêt agit le long de la ficelle, donc en passant
par l'axe~: aucun moment par rapport à lui, et le \omterm{def:b1:angular-momentum:L}{moment cinétique}
$J\omega$ n'est pas touché. Seule une force tangentielle (un frottement
sur l'axe) pourrait ralentir la rotation.

\textbf{15.} $T \approx Mv_G/\Delta t + Mg = 0.050 \times 0.82/0.005 + 0.49 =
\qty{8.7}{N}$, dix-huit fois le poids --- d'où la secousse que sent le doigt.

\textbf{16.} $J\dot\omega = -\Gamma$~: $t = J\omega/\Gamma = 2.25\times10^{-5} \times
205/10^{-4} = \qty{46}{s}$.

\textbf{17.} Toute l'énergie de rotation $\tfrac12 J\omega^2 = \qty{0.47}{J}$
se convertit en $Mgh'$~: $h' = \qty{0.96}{m}$, presque toute la hauteur~;
moins en pratique, à cause du frottement sur l'axe pendant le sommeil et
des pertes à la secousse.

\textbf{18.} $a_G$ est dirigée vers le bas (le yo-yo ralentit en montant)
alors que $v_G$ est vers le haut~: $T = M(g - a_G) < Mg$ --- plus petite
que le poids.

\textbf{19.} Pour $r = R$~: facteur $1 + 1/2 = 1.5$, $a_G = 2g/3 = \qty{6.5}{m/s^2}$
--- chute rapide, peu de rotation~: c'est l'axe fin qui fait le yo-yo.

\textbf{20.} \omterm{ex:b1:systems-of-points:rolling}{Roulement sans glissement}~: le \omterm{def:b1:point-kinematics:frame}{point matériel} en $C$ a une
vitesse nulle~; or le champ des vitesses d'un \omterm{def:b1:systems-of-points:solid}{solide} dont un point est
immobile est une rotation autour de ce point (axe passant par $C$, même $\omega$).

\textbf{21.} $F$ s'applique au bas de l'axe, à la hauteur $R - r$ au-dessus
de la table, horizontalement vers la main~: son moment vaut $F(R - r)$
par rapport à $C$, dans le sens qui roule la bobine vers la main~; le
poids et la réaction passent par $C$ (moment nul).

\textbf{22.} $(J + MR^2)\dot\omega = F(R - r)$, $a_G = R\dot\omega = FR(R - r)/
(\tfrac32 MR^2) = 2F(R - r)/3MR > 0$~: vers la main.

\textbf{23.} La droite d'action de $F$ passe par $C$ lorsque la ficelle,
partant du bas de l'axe, rencontre la table en $C$~: $\cos\beta = r/R$,
$\beta = 82^\circ$. Plus redressée~: le moment par rapport à $C$ s'inverse
et la bobine s'éloigne (et pour $\beta$ assez grand, elle se soulève).

\textbf{24.} $a_G = 2 \times 0.20 \times 0.026/(3 \times 0.050 \times 0.030) =
\qty{2.3}{m/s^2}$~; le frottement, par $Ma_G = F - f$ (il s'oppose à $F$)~:
$f = 0.20 - 0.050 \times 2.3 = \qty{0.085}{N}$~; il doit rester sous $\mu_sMg
\approx \qty{0.2}{N}$ (pour $\mu_s = 0.4$), c'est-à-dire $F \lesssim \qty{0.5}{N}$.

\textbf{25.} \omterm{def:b1:systems-of-points:cm}{Centre de masse}~: la translation de $G$ et la tension~;
\omterm{def:b1:angular-momentum:L}{moment cinétique} par rapport à un axe~: la rotation et le sens du
roulement~; énergie~: la vitesse en bas et la hauteur de la remontée.
\end{solution}
